% =============================================================================
% Section 2 -- Background and Related Work
% =============================================================================

\section{Background and Related Work}
\mohit{I would add more in the related work. It looks small. Also fix Table 1} 

\yash{this section is ready for review}


\candidate{S2.0}
When two states go to war, both publics argue about it online at the same time,
in communities aligned with one side, communities aligned with the other, and
general-audience communities where the conflict is argued at scale. One event
reaches all of them at once, and what each community makes of it is not the
same. That difference is what we measure.

Most of what such a measurement needs already exists, in pieces. Differently
aligned audiences are known to narrate the same conflict
differently~\cite{entman1991kal,makhortykh2017visual}. Harm is now treated as a
composite, with several labelled attributes of a comment modelled together
rather than one at a time~\cite{scheuerman2021severity,zhou2025harm}, and a
parallel line of research composes attributes of the person doing the
labelling~\cite{rottger2022paradigms,davani2024d3code}. Crossing between opposed
communities has been examined directly, and has generally been found to produce
friction~\cite{kumar2018conflict,bail2018exposure}. Norms themselves are
stratified rather than site-wide~\cite{chandrasekharan2018hiddenrules}, so the
same sentence is not the same object in two communities. What none of this
composes is the community a comment appeared in.

This section takes those three literatures in turn, on narrating conflict, on
measuring framing and harm, and on crossing between communities, and says where
each stops and why. Each stopping point is a property of an
instrument, by which we mean what a study measures with: what it labels, on what
unit, and from which recorded signal. An instrument bounds a study's findings,
because a quantity it never records cannot be recovered afterwards from the ones
it does.


% -----------------------------------------------------------------------------
\subsection{Narrating interstate conflict}
\label{sec:rw-conflict}
% -----------------------------------------------------------------------------

\subsubsection{Frames, and the inventories used to measure them}

\candidate{S2.1.1}
A frame, in the definition we adopt throughout, selects ``some aspects
of a perceived reality'' and makes them ``more salient in a communicating text,
in such a way as to promote a particular problem definition, causal
interpretation, moral evaluation, and/or treatment
recommendation''~\cite{entman1993framing}. The founding demonstration of that idea compared United States coverage of two
civilian airliner shootdowns~\cite{entman1991kal}. One was covered as moral
outrage, the other as technical accident, and the difference tracked whose side
had fired. Frame effects are not fixed, however. They depend on which competing
frames are available and on what the audience already
believes~\cite{chongdruckman2007framing}. That conditionality is the theoretical
reason to expect a frame to behave differently in two communities, and we test
it rather than assume it. Computational work measures frames
against inventories, and the established practice for a new domain is to take a
general-purpose inventory~\cite{card2015mediaframes} and add the frames the
domain requires~\cite{mendelsohn2021framing}. Our six frames follow that practice; Section~\ref{sec:data} gives their derivation.

\subsubsection{The same war, told by differently aligned audiences}

\candidate{S2.1.2}
The split Entman found recurs wherever a conflict has two audiences. Work on the war in eastern Ukraine found pro-Ukrainian communities framing it as
a limited action against insurgents, and pro-Russian communities framing the same fighting
as a war on civilians~\cite{makhortykh2017visual}. The pattern also appears
within one platform and one country. How \#MeToo was discussed depended on the
political orientation of the community discussing it~\cite{rho2018metoo}. Those
results establish that the pattern exists. The open question is what governs it.

One quantitative comparison bears directly on that question.
\citet{hanley2023specialoperation} compared three media ecosystems covering the
Russian invasion of Ukraine, and found Western outlets concentrating on its
military and humanitarian aspects, and Chinese outlets on its diplomatic and
economic consequences. Two things follow.
Aligned audiences covering one event divide along emphasis rather than accuracy,
which is the interaction we measure at comment level rather than outlet
level. And four of our six frames appear there as recurring emphases in an
independent corpus, which is evidence that our inventory was not built to fit
these data. The two it does not reach, historical grievance and religious or
cultural framing, are specific to this dyad, and
Section~\ref{sec:data} warrants them instead.

Other studies extend the observation in different directions.
\citet{mejova2025narratives} modelled Ukrainian wartime tweets as narratives with
assigned roles, and found that a victim narrative roughly doubled a tweet's
reach. On
a different dyad, \citet{shin2026diplomacy} separated framing from affect across a diplomatic
summit: sentiment reverted once talks failed, and the shift away
from threat-oriented framing did not. Work on the Rohingya crisis inverted the usual target and sampled supportive
speech rather than hostility~\cite{palakodety2020rohingya}. Within
CSCW, frame contests have been shown to be networked and observable in the
structure of interaction itself~\cite{stewart2017lines}, which established that a frame can be read off a community's structure rather
than only off its text. That work did not compare the structure across several
communities holding one event fixed, which is what we do instead. Affect under sustained armed
conflict has likewise been tracked over time~\cite{dechoudhury2014narco}.

\subsubsection{What these designs measure, and what they hold fixed}

\candidate{S2.1.3}
These studies share an instrument, and the instrument is what bounds them. Each
labels the speaker: the frame a post advances, the stance it takes, the
hostility it expresses. None labels what the receiving audience does with it.
That is not an oversight, since a corpus of posts supports audience-side
measures only where the platform records a reaction. But the quantity reported
here, whether a community divides over a comment, is outside the reach of these
designs rather than merely missing from them.

A second bound is narrower, and we only partly escape it. Each study
takes one conflict and usually one window of time. A pattern recovered from a
single window cannot afterwards be separated into the period it was recorded in
and the kind of event that produced it. We observe six crises across seven years, two of them military events six years
apart. That supports a
comparison between events of the same kind, which a single-window design cannot
run. It is a robustness check rather than a separate claim. Two of the three
crisis types here are single events, so nothing about crisis type can be
estimated at those levels, and Section~\ref{sec:data} states the
coding and its limits. The comparison rule is inherited:
\citet{olteanu2015crisis} compared discourse across crises by proportion within a
community rather than by raw volume.

\subsubsection{The India--Pakistan dyad, and what it supplies}

\candidate{S2.1.4}
Computational work on this conflict has been consistent in shape. Prior studies
have taken a single event, usually one crisis in 2019, labelled one or two
attributes per post, and worked on Twitter or YouTube~\cite{tyagi2020ipconflict,hussain2021pulwama}. A second programme has approached the dyad from the opposite direction,
detecting hope and peace-oriented speech rather than hostility, in the
code-mixed form the discourse is actually conducted in~\cite{palakodety2020hope,khudabukhsh2020codeswitch}. On Reddit, an adjacent
conflict has been labelled on one attribute, stance~\cite{ali2025stance}. \citet{chen2024isamasred} released a public dataset of roughly eight million
comments from Israel--Hamas discussion, carrying measures of controversy and
moral language, which is the outcome we also measure. They released those
measures as a dataset rather than composing them with frame or community, which
is the step we take next.

The setting supplies three things to the design rather than merely hosting it.
A community's alignment here follows the territorial dispute rather than any
elective affiliation, so the position belongs to the community itself rather
than to whichever speaker happens to post in it.
Participants are audiences to one shared event on a clock set outside the
platform, so comparing communities compares readings of one occurrence. And the
discourse already has a scholarly literature that a new measurement has to
answer to. Research in the
region has shown how gendered digital abuse is
organised~\cite{sambasivan2019alone}, and how religious minority communities in
Bangladesh experience political fear online~\cite{rifat2024politicsoffear}. Both bear
on the results rather than sitting beside them. Our six-frame inventory has no gendered axis, which bounds what it can see in a discourse partly organised
along one. The fear mechanism is one reading of the non-crossing reported in
Section~\ref{sec:data}. The gap this subsection leaves open is therefore a
divided one. Social computing has studied how harm is organised in this region,
and has studied how aligned audiences frame a shared conflict elsewhere, but has
not brought the two together. We are aware of no CSCW or HCI study that measures
frame-level discourse on this dyad, and the computational work that does measure
it has appeared largely outside those venues. We measure it inside them.

% -----------------------------------------------------------------------------
\subsection{Composing harm: what gets measured together, and what does not}
\label{sec:rw-measurement}
% -----------------------------------------------------------------------------

\subsubsection{One construct at a time, and the instrument that labels it}

\candidate{S2.2.1}
Harm is measured one construct at a time, because that is what annotation
budgets and pretrained classifiers afford. Offensive-language
benchmarks~\cite{zampieri2019olid} and resources for implicit
hate~\cite{elsherief2021latenthatred} are each built and validated alone. Where
their labels meet in a model, they meet as independent features. The training data underneath them are unevenly documented~\cite{vidgen2020garbage}, and
labeller error within it is distributed unevenly across
dialect~\cite{sap2019racialbias}. Fairness assumptions carried into a new locale
do not transfer intact~\cite{sambasivan2021india}. We measure our own comparison outcome against a score built this way, so we
report our offensiveness finding as a result about that particular instrument,
not about offensiveness in general.
Section~\ref{sec:data} names the instrument and its threshold.

The instrument that does the labelling has also changed. Prompting a
general-purpose language model against a written codebook is now ordinary
practice for social-scientific
constructs~\cite{gilardi2023chatgpt,ziems2024llmcss}, with performance
contingent on the dataset and the task~\cite{pangakis2023validation}, and a
standards literature specifies what such a pipeline must
report~\cite{tornberg2024standards}. That literature assumes a reliable human
reference. For contested political constructs the human labels themselves agree
only moderately, and across thousands of annotation tasks the best models sit
close to human--human agreement rather than far below
it~\cite{kunilovskaya2026annotators}, so that level of agreement is the field's
normal condition rather than a local failure. The design this affords is
established in social computing: hand-code a sample, train or prompt a labeller
to reproduce that code on the rest of the corpus, then compare rates across
communities, as in a comparison of muting across seven
subreddits~\cite{foriest2024muting} and of frames across an aligned
public~\cite{panda2026framing}. Both used supervised classifiers rather than prompted language models, so what
we inherit is the design and not the instrument; Section~\ref{sec:data} reports what its labels agree with.

\subsubsection{Composition, and what it composes}

\candidate{S2.2.2}
Composition already exists, and it is CSCW's own work. \citet{scheuerman2021severity} developed a severity framework from interviews and
card-sorting, setting out four types of harm and eight dimensions, and argued
that the dimensions may compound on one another. \citet{zhou2025harm} composed
three attributes over a corpus of anti-Asian hate speech: the presence of
misinformation, the collective narrative, and the targeted entity. We extend both rather than correct either, and saying where they stop requires
precision. Only one of the severity framework's eight dimensions is contextual
in the sense needed here: sphere, which separates public posts from private
messages, a coarser distinction than the one between one community and another.
\citet{zhou2025harm} composed attributes of the utterance, and their second
research question added factors belonging to the person judging. That varies who is doing the judging, the
perspectivist axis, rather than where the comment was posted. Their corpus is
Twitter, and the subreddit does not appear in
it. Their outcome also differs in kind: they measured harmfulness as judged by a
reader, where we measure whether a community divided.

A second line composes the coder rather than the content. It treats
disagreement between annotators as signal rather than noise, modelling who
labels instead of averaging over
them~\cite{rottger2022paradigms,sap2022annotators,davani2024d3code}, and systems
work has begun to make that disagreement something to deliberate over rather
than resolve, including by placing models inside the deliberation
itself~\cite{khadar2025socratic}. That programme varies the coder. It does not
ask whether the same comment is a different object in a different community,
which is an adjacent question rather than a criticism of it. Between the two
lines, the attributes of a comment and the attributes of the person coding it
have both been composed. Where the comment was posted has not.

\subsubsection{The community as a measured object}

\candidate{S2.2.3}
Norms on Reddit are stratified. \citet{chandrasekharan2018hiddenrules} studied moderator removals across a
hundred of the site's largest communities and found three levels. Some norms are enforced across
most of the platform, some hold across groups of communities, and some are
specific to one. What became of that finding is instructive. Crossmod, built in
the same programme, scores a comment with an ensemble of one hundred and eight
classifiers~\cite{chandrasekharan2019crossmod}: one hundred predict whether one
subreddit's moderators would remove it, and eight cover site-wide norms. The
tool operationalises the site-wide and single-community layers and has no
counterpart for the intermediate one. A stratification discovered in the data
was only partly carried into the instrument built from it. Our alignment groups are an attempt at that intermediate layer, assigned rather
than learned.

% -----------------------------------------------------------------------------
\subsection{What communities reward, and what happens at their boundary}
\label{sec:rw-crossing}
% -----------------------------------------------------------------------------

\subsubsection{Contestation, and how it differs from hostility}

\candidate{S2.3.1}
We compose narrative frame with community to explain what makes crisis discourse
contested, and that composition tells us something new only if contestation is
not just another name for hostility. If a community divides over
a comment exactly when the comment is hostile, then a toxicity score, which harm
research already measures, would answer the same question, and community would
add nothing beyond a variable to control for. So the outcome we measure is contestation: whether a community divided over a comment, not whether the
comment was hostile. The two come apart in principle, and Section~\ref{sec:data}
reports that they come apart in these data, which is what licenses treating
contestation, rather than the toxicity score, as the outcome composed with frame
and community below.

Measuring contestation this way has two precedents. Controversy has been
treated as a structural property of interaction, readable from how a discussion
network splits~\cite{garimella2016controversy,garimella2018controversy}. It has
also been predicted from the early shape of a discussion, using the signal
Reddit itself records~\cite{hessel2019controversy}. That signal, whether a comment drew votes both ways, is the one we use.
Section~\ref{sec:data} states its limits: it registers votes falling
both ways rather than the quality of the exchange, and requires a comment to
attract enough votes to be assessed at all.

Reading a vote this way, as a community dividing rather than merely reacting,
requires a further warrant: that what a community rewards is a property of the
community and not an artefact of whoever happened to comment that day.
Communities differ systematically in what they reward as well as in what they
say~\cite{waller2021quantifying}, and can state what they value in terms
recognisable to themselves~\cite{weld2024values}. Affective polarisation and
partisan sorting are the broader conditions these communities sit
inside~\cite{iyengar2019affpol}. Newcomers acquire
the community-specific norm layer quickly~\cite{rajadesingan2020norms}, so a
reward pattern can be read as a property of the community rather than an
aggregate of whoever arrived. That is the premise we need to compose contestation with community in the first
place, rather than treating the
community a comment appeared in as background.

\begin{table*}[t]
\footnotesize
\caption{Computational work on this dyad and on the nearest adjacent cases.}
\label{tab:positioning}
\newcolumntype{Y}[1]{>{\hsize=#1\hsize\raggedright\arraybackslash}X}
\begin{tabularx}{\linewidth}{@{}Y{1.15}Y{0.70}Y{1.05}Y{1.30}Y{1.05}Y{1.03}Y{0.72}@{}}
\toprule
\textbf{Study} & \textbf{Platform} & \textbf{Events} & \textbf{Labels and instrument} & \textbf{Communities} & \textbf{Community in the measurement} & \textbf{Language} \\
\midrule
Tyagi et al.~\cite{tyagi2020ipconflict} & Twitter & Pulwama, Balakot & Hashtag labels & Politicians & Not modelled & English \\
Hussain~\cite{hussain2021pulwama} & Twitter & Pulwama & Manual, four codes & Two national sides & Not modelled & English \\
KhudaBukhsh et al.~\cite{khudabukhsh2020codeswitch} & YouTube & 2019 crises & Code-switching, classifier & Not modelled & Not modelled & Code-mixed \\
Palakodety et al.~\cite{palakodety2020hope} & YouTube & Pulwama, Balakot & Hope speech, classifier & Not modelled & Not modelled & Code-mixed \\
Ali et al.~\cite{ali2025stance} & Reddit & Israel--Hamas & Stance, one layer & Two & Grouping variable & English \\
Chen et al.~\cite{chen2024isamasred} & Reddit & Israel--Hamas, one window & Controversy, moral language & Not modelled & Not modelled & English \\
Rao et al.~\cite{rao2025abortion} & Twitter & Abortion, one year & Five frames, hostility; transformers & Liberal, conservative (per user) & Speaker attribute & English \\
Panda et al.~\cite{panda2026framing} & Twitter & One campaign & Grounded-theory frames; RoBERTa & One aligned public & Grouping variable & English \\
Murad Ahmed et al.~\cite{muradahmed2026ukraine} & Twitter & Ukraine 2014 & Themes, hashtag projection & Three blocs incl. neutral & Grouping variable & English \\
Zhou et al.~\cite{zhou2025harm} & Twitter & COVID-19 period & Misinformation, narrative, target & Anti-Asian hate & Not modelled & English \\
\midrule
\textbf{Our study} & \textbf{Reddit} & \textbf{Six, 2019--2025} & \textbf{Six labelled constructs, LLM against a codebook, human-validated} & \textbf{Ten, in three alignment groups} & \textbf{Measured: composed with frame, and the boundary crossed} & \textbf{Hindi, Urdu, English} \\
\bottomrule
\end{tabularx}
\end{table*}



\subsubsection{Crossing, and the instrument that identifies a bridge-user}

\candidate{S2.3.2}
What happens at the boundary a participant crosses, from a community aligned
with one side of the conflict into one aligned with the other, is the third question our design has to answer, and the existing literature answers
it only provisionally. Where cross-community participation has been examined
directly, exposure to the
other side has generally produced friction rather than
accommodation~\cite{kumar2018conflict,bail2018exposure,datta2019interconflict}. That
result is usually reported against intergroup contact theory, whose
meta-analytic statement makes reduced prejudice conditional on facilitating
conditions~\cite{pettigrew2006contact}: equal status, common goals and
institutional support. None of these conditions typically holds when a participant posts into a
community aligned with the other side of a shooting war. Read that way, friction is the
predicted case rather than a surprising one.

The picture is not unanimous, and the disagreement matters for our argument
rather than against it. Behaviour after cross-community contact depends on
conditions~\cite{xia2025integrated}, cross-partisan discussion can be better sustained under particular
designs~\cite{rajadesingan2021crosspartisan}, and
antisocial behaviour is at least partly situational rather than
dispositional~\cite{cheng2017troll}. Silence is itself a mechanism rather than
an absence of speech~\cite{noelleneumann1974spiral}, which is one reading of
the non-crossing figures in Section~\ref{sec:data}.

Two measured precedents share our subject but not our instrument.
\citet{wu2021crosspartisan} showed on YouTube that crossing is asymmetric,
conservatives being far more likely to comment on left-leaning videos than the
reverse, and that cross-partisan replies are more toxic than co-partisan ones.
\citet{almerekhi2022toxicity} showed at very large scale that most users who
comment in more than one community changed their toxicity across them. Both
measure the crosser's own hostility; neither measures what the receiving
community does, which is the quantity separating a penalty on an outsider from a
change in an outsider's writing. Community-level estimates are also a poor guide
to particular participants: \citet{trujillo2022reddit} found a moderation
intervention falling differently on core and peripheral users of the same
community, which warrants following individuals rather than aggregates.

That instrument bounds what this literature can say in two ways. The first is
how a crosser is identified. Social computing has located them in the
structure around the person: in network position, where inter-community conflict
is measured as mobilisation between groups rather than as the movement of any
one participant~\cite{kumar2018conflict,datta2019interconflict}; in how a user's
activity distributes across communities~\cite{waller2019generalists}; or in the
hostility their comments carry in
aggregate~\cite{wu2021crosspartisan,almerekhi2022toxicity}. Each of these
conflates two things that move together: being legible as an outsider there, and
writing differently there, so a penalty measured this way cannot be assigned to
either cause. The second bound is that a crosser is treated as a kind of person
rather than as a matter of
degree. Activity diversity is continuous, and generalist and specialist are ends
of one spectrum rather than two populations~\cite{waller2019generalists}, so any threshold sorting users into bridges and non-bridges imposes a division
the data do not contain.

Together, these bound what CSCW can do with the result, because the field's
interest in cross-community contact is finally a design interest: the question is what platform designers or a community's own moderators could
change~\cite{rajadesingan2021crosspartisan}. That
question has different answers depending on which penalty operates. A community
penalising an outsider's alignment would not respond to an intervention aimed at
how people write, and one penalising their register, the tone and vocabulary a
comment is written in rather than who wrote it, might. Deciding between them
needs an instrument that holds the person fixed, which is what our within-author comparison does: it compares the comments an author posts inside their
own community against the ones the same author posts across the boundary, so
identity is constant and register is free to vary.


% -----------------------------------------------------------------------------
\subsection{Positioning}
\label{sec:rw-positioning}
% -----------------------------------------------------------------------------

\candidate{S2.4}
We extend each of these literatures. From work on how aligned audiences narrate
a shared conflict we take the frame-by-audience interaction, moving it from
outlets to communities and from one event to a panel, several occurrences of the
same kind of crisis observed over time. From the harm-measurement literature we
take composition, which CSCW has already established, and add the community to
what is composed. The community becomes
part of what is measured rather than a variable held constant. From the boundary-crossing literature we take the friction
finding, and supply the within-author comparison separating a penalty on a
crosser's alignment from a penalty on their register.

\yash{add about author's background here? or in the data annotation section?}


